% Nouveau EGA V (cote 25): archive cote 25, site page 148 (PDF page 149).
% typescript. Draft compared: none.
% AI-transcribed by Claude Opus 5.5 (claude-opus-5-5, adaptive thinking, effort high), Message Batches API, batch ega25b (msgbatch_01A7n5YuXubUPNK1zsepi8RG), 2026-11-13; not yet checked by eye.
\scanpage{25/148}
\mspage{\hand{9}}

ment normal de \var{\xout{X}}{}{typewriter strike-out}{ega25_25_148-v0001}$V_y$ sur k(y), (puisque de toutes façons \var{\xout{X}}{}{typewriter strike-out}{ega25_25_148-v0002}  $V_y$ est CM ).
On peut montrer  plus généralement, sous les conditions de déf.1 , que si
x est un zéro singulier ordinaire, alors x est une singularité isolée
de V($\phi$)  en ce sens que toute générisation \var{\xout{\unsure{de}}}{}{typewriter strike-out}{ega25_25_148-v0003} x' de x dans V($\phi$) est un
point régulier de V($\phi$)  (donc si le lieu singulier de V($\phi$) est fermé, p.
ex. si X est ``ex\var{\xout{c}}{}{typewriter strike-out}{ega25_25_148-v0004}cellent'' , x est \var{\xout{bien}}{}{typewriter strike-out}{ega25_25_148-v0005} un point isolé de ce lieu singuli-%
er). \var{\xout{Lorsque}}{}{typewriter strike-out}{ega25_25_148-v0006} (Pour bien faire, il faudrait inclure ce résultat au Chap
IV , mais  le par. 16 ne semble guère le lieu. Où le mettre ?).  Il
s'ensuit que si X est lisse sur un corps k, et si x est un zéro géométri%
quement singulier ordinaire, alors  x est un point isolé de l'ensemble
des points de non lissité  de V(\var{\over{\unsure{4}}{$\phi$}}{$\phi$}{writing over}{ega25_25_148-v0007}) , donc est \var{\inserted{(géom.)}}{(géom.)}{insertion}{ega25_25_148-v0008} normal si  $\dim_x(X) \geqslant 2$ ,
(mais \var{\inserted[typed]{nécessairement}}{nécessairement}{insertion}{ega25_25_148-v0009} non normal si $\dim_x X = 1$ , comme on le \var{\deleted{vérifie}}{}{deletion}{ega25_25_148-v0010} \var{\inserted{constate}}{constate}{insertion}{ega25_25_148-v0011} par exemple sur
le cas type $\phi = xy$ , x et y fonctions coordonnées dans le plan affine).

\hand{Voici le diagramme analogue à (D), mais pour}
\hand{une section $\varphi$ d'un Module inversible K :}

\hand{(D')}
\diagram{
\begin{tikzcd}[ampersand replacement=\&,nodes in empty cells]
 \& \& 0 \& \& \\
 \& \& \Omega^1_{X/Y}\otimes\underline{O}_{V'} \arrow[u] \& \& \\
0 \arrow[r] \& P^1_{X/Y}(K)^{\vee}\otimes\underline{O}_{V'} \arrow[r] \arrow[ur, "\mu"] \& \Omega^1_{X/S}\otimes\underline{O}_{V'} \arrow[r] \arrow[u] \& \Omega^1_{V'/S} \arrow[r] \& 0 \\
 \& K^{-1}\otimes\underline{O}_{V'} \arrow[u] \arrow[r] \& \Omega^1_{Y/S}\otimes\underline{O}_{V'} \arrow[u] \arrow[ur, "\nu"'] \& \& \\
\omega^{\unsure{-2}}_{X/Y}\otimes K^{\unsure{-2n-1}}\otimes\underline{O}_{V'} \arrow[ur] \& \& 0 \arrow[u] \& \&
\end{tikzcd}
}

\mspage{\hand{148}}
